\section{历史回顾：LLM 之前的智能体方法}

在大语言模型出现之前，研究者已经探索了三种构建指令跟随智能体的主要方法。

\subsection{方法一：语义解析（类机器翻译）}

\begin{figure}[H]
\centering
\includegraphics[width=0.95\textwidth]{slides-images/slide-045.jpg}
\caption{方法一：将自然语言指令直接映射为动作序列，类似机器翻译。适用于简单的 text-to-SQL、知识图谱查询等场景\protect\footnotemark}
\end{figure}
\footnotetext{Slides 第46页。Source: Zettlemoyer et al. 2012。}

核心思路：将自然语言指令视为``源语言''，将可执行的逻辑表示视为``目标语言''，直接训练一个序列到序列模型来完成映射：

$$\max_\theta p_\theta(\{a_1, a_2, \ldots\} \mid g)$$

适用于简单的接地环境如 text-to-SQL 和知识图谱查询。

\subsection{方法二：推断计划 + 执行模型}

不直接映射指令到动作，而是：
\begin{enumerate}
    \item 从（指令，动作轨迹）对中推断出\textbf{抽象计划}
    \item 训练模型将指令映射到计划
    \item 定义执行模型来执行这些计划
\end{enumerate}

优势在于计划可以编码更高层次的决策，而不仅仅是低层动作序列。

\subsection{方法三：强化学习}

直接使用 RL 学习以自然语言指令和环境观测为条件的策略：

$$\pi(a_t \mid \text{observation}_t, g)$$

奖励可以是稀疏的（完成整个任务后才获得反馈）或密集的（每一步都有反馈）。

\begin{knowledgebox}{三种方法的对比}
\begin{itemize}
    \item \textbf{语义解析}：需要大量标注的（指令，逻辑形式）对，但对简单领域效果好
    \item \textbf{计划推断}：可以处理更复杂的任务，但需要设计执行模型
    \item \textbf{强化学习}：最通用，但需要环境反馈和大量交互
\end{itemize}
这些方法都需要\textbf{针对特定环境}从头训练模型，缺乏泛化能力。
\end{knowledgebox}

\subsection{本章小结}

LLM 之前的智能体方法虽然在各自的适用场景内有效，但都需要为每个新环境重新训练模型。LLM 的出现带来了新的可能性：利用预训练中获得的广泛知识，以最少的适配来处理各种环境中的任务。

%% ========== Part 7: LLM 智能体 ==========

